\section{A sharp shallow Gaussian oracle on arbitrary input laws}
\label{sec:general-shallow-sharp}
This section uses the Gaussian model, normalized path mass $P$, and
empirical midpoint inequality (GM1) of the preceding section. Inputs
have any probability law on $[-1,1]^d$; initially $|Y|\le1$.
There is no density, smoothness, or input-separation assumption.
All constants below are numerical. Put
\[
 b_n=\frac{c_0\log^4(\mathrm e(n+1)(d+1)/\delta)}n,
 \qquad \lambda_n=c_1\{\sqrt{b_n}+(d+2)b_n\},                 \tag{GS1}
\]
where $c_0,c_1$ are sufficiently large fixed numerical constants.
The nontrivial finite-sample theorem assumes $\lambda_n\le1$.
For each fixed $d,\delta$, this holds for all sufficiently large $n$.
The likelihood schedule knows neither the comparison coefficients nor
its path mass or risk. Its physical observation variance remains
$n\lambda_n\nu/N$.

\begin{lemma}[Dimension-free atomic and multiplier bounds]
\label{lem:general-shallow-statistics}
Let
\[
 \mathcal H_d=\{x\mapsto d^{-1/2}(\theta\cdot\bar x)_+:
                              \|\theta\|_2=1\}.
\]
Its elements are bounded by $\sqrt2$, and the absolute convex hull
with coefficient mass at most $p$ has empirical Rademacher complexity
at most $Cp/\sqrt n$, for every input sample.
For any deterministic finite shallow ReLU function $g$ satisfying
$R(g)\le2$, with probability at least $1-\delta/2$,
\[
 \sup_{u\in\mathcal H_d}
 |(\mathbb P_n-\mathbb P)[(g(X)-Y)u(X)]|
 \le C\left\{\sqrt{\frac{\log(4/\delta)}n}
       +(P_g+1)\frac{\log(4/\delta)}n\right\}.              \tag{GS2}
\]
Here $\mathbb P$ denotes the joint input-label law; $P_g$ is path
mass and is unrelated to the probability notation.
\end{lemma}
\begin{proof}
ReLU contraction and $\|\bar x\|_2/\sqrt d\le\sqrt2$ give
\[
 \frac1n\mathbb E_\epsilon\sup_{\|\theta\|=1}
 \left|\sum_i\epsilon_i d^{-1/2}(\theta\cdot\bar X_i)_+\right|
 \le\frac Cn\mathbb E_\epsilon
           \left\|\sum_i\epsilon_i\bar X_i/\sqrt d\right\|_2
 \le C/\sqrt n.
\]
Absolute convexification multiplies this by $p$.

For (GS2), write $a_i=g(X_i)-Y_i$. The weighted contraction inequality
has the same proof, since each scalar map
$s\mapsto a_i(s/a_i)_+$, for $a_i\ne0$, is $1$-Lipschitz and vanishes
at zero; zero $a_i$ contributes zero. Conditional on the data, its
bound is $C\sqrt{\sum_i a_i^2}/n$. Symmetrization and Jensen therefore
bound the expected supremum of the centered empirical process by
$C\sqrt{R(g)/n}\le C/\sqrt n$.
Each summand has variance at most $2R(g)\le4$ and absolute envelope
at most $C(P_g+1)$, because $|g|\le\sqrt2P_g$.
The bounded empirical-process Bernstein inequality (Bousquet's
inequality~\cite{bousquet}) gives, with $t=\log(4/\delta)$, the bound
\[
 C\left[n^{-1/2}+\sqrt{t/n}
       +\sqrt{(P_g+1)t/n^{3/2}}+(P_g+1)t/n\right].
\]
The mixed square-root term is at most a constant times
$\sqrt{t/n}+(P_g+1)t/n$, proving (GS2).
Continuous parameterization on the compact sphere supplies a
countable determining subclass; no measurability issue arises.
\end{proof}

\begin{lemma}[A relative square-norm event]
With probability at least $1-\delta/2$, uniformly over all finite
shallow ReLU functions $h$ of path mass $P_h\le 4n^5$,
\[
 \mathbb E h(X)^2-\mathbb P_n h(X)^2
 \le\tfrac14\mathbb P_n h(X)^2+C b_n(P_h+1)^2.              \tag{GS3}
\]
\end{lemma}
\begin{proof}
Apply the smooth nonnegative-loss Rademacher bound of Srebro,
Sridharan, and Tewari~\cite[Theorem~1]{sst} to square loss with zero labels on path balls
of radii $1,2,4,\ldots,2^{\lceil\log_2(4n^5)\rceil}$.
On a radius-$p$ ball the loss is bounded by $2p^2$, is $2$-smooth,
and the preceding lemma gives worst-case Rademacher complexity
$Cp/\sqrt n$. Allocate confidence equally among these
$O(\log n)$ balls. The resulting upper deviation is bounded by
\[
 C\sqrt{\mathbb P_nh^2}\,p\sqrt{b_n}+C p^2b_n.
\]
Young's inequality and the least dyadic radius containing $h$ prove
(GS3). This event depends only on a function class, so it is common
to every learner width and Gaussian scale.
\end{proof}

\begin{lemma}[A uniform posterior energy cutoff]
For $N\ge n$, $0<\nu\le1$, $n^{-2}\le\lambda\le1$, there is a
numerical constant $A$ such that, for
\[
 q_*=A\{d+2+\lambda^{-1}\},
\]
every sample with $|Y_i|\le1$ satisfies
\[
 \Pi_\lambda[(1+Q^2)\mathbf1_{\{Q>q_*\}}]\le n^{-40}.     \tag{GS4}
\]
\end{lemma}
\begin{proof}
The prior is centered and has pointwise output variance at most
$\nu^2/N$, so its empirical risk is at most two. Consequently
$Z\ge\exp(-\omega/\lambda)$. With $D=(d+2)N$,
$\mathbb E_\mu e^{\omega Q/4}=2^{D/2}$.
The elementary bound
\[
 (1+Q^2)\mathbf1_{Q>q}
 \le C(1+\omega^{-2})e^{\omega Q/4-\omega q/8}
\]
therefore gives
\[
 \Pi_\lambda[(1+Q^2)\mathbf1_{Q>q}]
 \le C(1+\omega^{-2})
       \exp\{\omega/\lambda+D\log2/2-\omega q/8\}.
\]
Use $D/\omega=\nu(d+2)\le d+2$, $\omega\ge n$, and enlarge $A$.
The exponent is then sufficiently negative to imply (GS4),
uniformly in all the stated parameters.
\end{proof}

\begin{theorem}[A coefficient-one path oracle]
\label{thm:general-shallow-sharp}
Suppose (GS1) satisfies $\lambda_n\le1$. For any deterministic finite
shallow ReLU comparator $g$, at every $N\ge n$ and $0<\nu\le1$,
with probability at least $1-\delta$,
\begin{align}
 R(\bar z_{\lambda_n})\le\Pi_{\lambda_n}R(z)
 \le R(g)+C\lambda_n\left[
 P_g+1+(d+2)\log((n+1)(d+1))\right].                     \tag{GS5}
\end{align}
The event can be chosen common to all the eligible learner widths
and Gaussian scales. The prior and prediction are unchanged; the
likelihood scale is deterministic and independent of the comparator.
\end{theorem}
\begin{proof}
First assume $R(g)\le2$ and $\lambda_n(P_g+1)\le c$ for a sufficiently
small fixed constant $c$. In particular $P_g\le n^4$.
Intersect (GS2) and (GS3), and write $h=z-g$ and $P=P(z)$.
Since $h$ has path mass at most $P+P_g$, expand the squared-loss
difference and use these events to obtain
\begin{align*}
 R(z)-R(g)\le{}&\widehat R(z)-\widehat R(g)
       +\tfrac14\mathbb P_n(z-g)^2+C b_n(P+P_g+1)^2\\
       &+C(P+P_g)\{\sqrt{b_n}+(P_g+1)b_n\}.
\end{align*}
For $Q\le q_*$, $P\le q_*/2$. The chosen scale satisfies
\[
 b_nq_*\le c'\lambda_n,\qquad
 \sqrt{b_n}+(P_g+1)b_n\le c'\lambda_n,
\]
with $c'$ as small as needed by increasing $c_1$.
Indeed $\lambda_n\ge c_1\sqrt{b_n}$ and
$\lambda_n\ge c_1(d+2)b_n$; also
$b_n(P_g+1)\le c b_n/\lambda_n$.
The comparator-only terms obey
$b_n(P_g+1)^2\le C\lambda_n(P_g+1)$.
Hence on this energy set
\begin{align*}
 R(z)-R(g)\le{}&\widehat R(z)-\widehat R(g)
       +\tfrac14\mathbb P_n(z-g)^2
       +\tfrac12\lambda_n P+C\lambda_n(P_g+1).
\end{align*}
All path masses here are at most $4n^5$ after increasing a harmless
minimum sample constant; $\lambda_n\le1$ implies $d\le Cn$.

Integrate on $\{Q\le q_*\}$ and use (GM1).
Replacing the restricted empirical integral by its full integral
costs at most $\widehat R(g)\Pi(Q>q_*)$ because the other summands
are nonnegative. The population tail costs at most
$C\Pi[(1+Q^2)\mathbf1_{Q>q_*}]$.
Here $\widehat R(g)\le C(P_g+1)^2$ is polynomial in $n$; (GS4)
makes both corrections negligible. Discard the remaining positive
empirical-distance and path terms in (GM1). This proves (GS5).

For completeness, apply exactly the same proof to $g=0$ (without
requiring $\lambda_n\le c$): its multiplier has bounded envelope,
and all constants are absorbed by (GS1). This gives
\[
 \Pi R(z)\le R(0)+C\lambda_n[1+(d+2)\log((n+1)(d+1))].
\]
If $R(g)>R(0)$ this already implies (GS5).
If $\lambda_n(P_g+1)>c$, use $R(0)\le1$ and enlarge the constant
on $\lambda_n(P_g+1)$ to dominate the possible difference between
$R(0)$ and $R(g)$. This handles all remaining comparators, including
ones with $P_g>n^4$. Jensen gives the posterior-mean conclusion.
\end{proof}

\paragraph{Width and coefficient budgets.}
For an ordinary affine-output shallow teacher of width at most $C$
whose coefficients and biases have magnitude at most $B\ge1$, one
constant neuron implements its output bias and
$P_g\le2d(C+1)B^2$. Thus at fixed $d,B$, (GS5) has remainder
$\widetilde O((C+1)/\sqrt n)$, with no density assumption and no
cap on teacher or learner outputs. More explicitly the remainder is
\[
 \widetilde O\left(
 \left[n^{-1/2}+\frac d n\right]
                  [d(C+1)B^2+d]\right).                 \tag{GS6}
\]
The algorithm does not use $C$ or $B$. For every fixed declared
budget class, choose a deterministic population-risk minimizer before
applying (GS5); one exists by continuity on its compact finite
parameter set. Alternatively, a single deterministic $n^{-10}$-near
minimizer suffices, with its tolerance absorbed in the remainder.
For simultaneous integer budgets, keep the same scale (GS1), apply
the multiplier event to the at most $n^4$ pairs $1\le C,B\le n^2$
with failure allocation $\delta/[2(n+1)^4]$, and enlarge the fixed
constants already used in $b_n$. This changes only its logarithmic
constant, not the scale between classes. Larger budgets have
$\lambda_n d(C+1)B^2\ge c$, so the zero-comparator event gives their
stated resource bound directly. No prior over these classes is used.
